\documentclass{article}
\usepackage[utf8]{inputenc}
\usepackage{xcolor}
\usepackage{amsmath}
\usepackage{graphicx}
\usepackage{float}
\usepackage{minted}
\usepackage{circuitikz}
\usepackage{tikz}
\usetikzlibrary{shapes, arrows, calc, positioning, circuits.logic.US, circuits.logic.IEC}
\usepackage{geometry}
\geometry{a4paper, margin=1in}
\usepackage{hyperref}
\usemintedstyle{trac}

\usepackage{fancyhdr}
\pagestyle{fancy}
\fancyhf{}
\renewcommand{\headrulewidth}{0pt}
\fancyfoot[C]{\small 010153101 Digital Circuit and Microprocessor Fundamental}
\fancyfoot[R]{\thepage}

% Redefine plain style to match fancy style (for the title page)
\fancypagestyle{plain}{%
  \fancyhf{}%
  \fancyfoot[C]{\small 010153101 Digital Circuit and Microprocessor Fundamental}%
  \fancyfoot[R]{\thepage}%
  \renewcommand{\headrulewidth}{0pt}%
}

\title{Laboratory Exercise 1 \\ Logic Gates and Boolean Algebra}
\author{}
\date{}

\begin{document}

\maketitle

\section*{Introduction}
This laboratory exercise accompanies Lecture 2 (Logic Circuits). You will confirm De Morgan's theorem with real TTL integrated circuits, and then go from a word problem to a truth table, a Boolean expression, a schematic, a simulation, and finally a working circuit on the DE0-CV board. No VHDL is needed in this lab: all circuits are built on the breadboard or drawn as schematics in Quartus Prime.

\section*{Objectives}
\begin{enumerate}
    \item To build a circuit from NAND, NOT and OR gates and to measure its logic levels.
    \item To verify De Morgan's theorem experimentally.
    \item To derive a Boolean expression in Sum-of-Products (SOP) form from a truth table.
    \item To use Boolean algebra to simplify an expression and to prove by simulation that the simplified circuit is equivalent to the original.
    \item To draw a schematic, simulate it with a waveform file and implement it on the DE0-CV board.
\end{enumerate}

\section*{Equipment Required}
\begin{itemize}
    \item DC Power Supply (+5V), Digital Multimeter
    \item Breadboard and Jumper Wires
    \item TTL ICs: 74LS00 (NAND), 74LS04 (NOT), 74LS32 (OR)
    \item Resistors: $330\Omega$ (for LEDs and pull-down)
    \item 4-position DIP switch
    \item LEDs (Red/Green)
    \item Intel DE0-CV board and a PC with Quartus Prime
    \item (Optional) Logisim or Logisim-evolution: a light-weight tool for drawing and simulating small circuits at home or in the laboratory. It cannot program the DE0-CV board, so Quartus Prime is still needed for the board.
\end{itemize}

\section*{Pre-lab (complete before the session)}
The laboratory session is two hours long, so the paper work is done in advance.
\begin{enumerate}
    \item Read the \textbf{switch input circuit} and the \textbf{LED indicator stage} of Lab 0. Every gate output that you observe on the breadboard is connected to an LED indicator stage (an inverter of the 74LS04 followed by an LED and a $330\,\Omega$ resistor to $+5\,\mathrm{V}$), so the LED is ON when the observed output is logic 1. The pin numbers of the ICs of this lab are given in Table \ref{tab:pinout}.
    \item \textbf{Tool tip.} You can check your pre-lab results at home with \textbf{Logisim}: draw the circuit with Input pins, gates and an Output pin and toggle the inputs with the Poke tool; the \emph{Combinational Analysis} window (Window menu) turns a truth table into an expression and a circuit, and shows the minimized expression.
    \item \textbf{Three-way light.} A room has three doors, each with a switch ($x_1, x_2, x_3$) that controls one light $f$. The light must toggle whenever any switch is flipped, and it is OFF when all switches are open. Complete the truth table (Table \ref{tab:three_way}) and write the canonical SOP expression: \\[0.2cm]
    $f = $ \underline{\hspace{10cm}}
    \item \textbf{Car safety alarm.} The alarm $A$ must sound if (a) the key is in ($K=1$) and the door is not closed ($D=0$), \textbf{or} (b) the door is closed ($D=1$), the key is in ($K=1$), the driver is in the seat ($S=1$) and the seat belt is not fastened ($B=0$). \\[0.2cm]
    Expression from the words: $A_1 = $ \underline{\hspace{7cm}} \\[0.2cm]
    Simplified with Boolean algebra (show the steps and the properties used; hint: factor $K$, then use $X + X'Y = X + Y$): \\[0.2cm]
    $A_2 = $ \underline{\hspace{7cm}}
\end{enumerate}

\begin{table}[H]
    \centering
    \begin{tabular}{|l|c|c|c|c|c|c|}
        \hline
        \textbf{IC} & \textbf{Gate 1} & \textbf{Gate 2} & \textbf{Gate 3} & \textbf{Gate 4} & \textbf{GND} & \textbf{$V_{CC}$} \\
        \hline
        74LS00 (NAND), 74LS32 (OR) & 1, 2 $\rightarrow$ 3 & 4, 5 $\rightarrow$ 6 & 9, 10 $\rightarrow$ 8 & 12, 13 $\rightarrow$ 11 & 7 & 14 \\
        \hline
        74LS04 (NOT) & \multicolumn{4}{c|}{1$\rightarrow$2, 3$\rightarrow$4, 5$\rightarrow$6, 9$\rightarrow$8, 11$\rightarrow$10, 13$\rightarrow$12} & 7 & 14 \\
        \hline
    \end{tabular}
    \caption{Pin numbers of the TTL ICs used in this laboratory. Each gate is listed as (input, input $\rightarrow$ output).}
    \label{tab:pinout}
\end{table}

\begin{table}[H]
    \centering
    \begin{tabular}{|c|c|c|c|c|}
        \hline
        \textbf{Minterm} & $x_1$ & $x_2$ & $x_3$ & $f$ \\
        \hline
        $m_0$ & 0 & 0 & 0 & \\
        $m_1$ & 0 & 0 & 1 & \\
        $m_2$ & 0 & 1 & 0 & \\
        $m_3$ & 0 & 1 & 1 & \\
        $m_4$ & 1 & 0 & 0 & \\
        $m_5$ & 1 & 0 & 1 & \\
        $m_6$ & 1 & 1 & 0 & \\
        $m_7$ & 1 & 1 & 1 & \\
        \hline
    \end{tabular}
    \caption{Truth table of the three-way light control.}
    \label{tab:three_way}
\end{table}

\section*{Part I: De Morgan's Theorem on the Breadboard (about 30 min)}

De Morgan's theorem states that $(A\cdot B)' = A' + B'$ (and $(A+B)' = A'\cdot B'$). You will check the first law with real gates.

\begin{figure}[H]
    \centering
    \begin{circuitikz}[circuitikz/logic ports=ieee]
        % Upper circuit: NAND
        \draw (0,2) node[nand port] (nand) {};
        \draw (nand.in 1) -- ++(-1,0) node[left] {$A$};
        \draw (nand.in 2) -- ++(-1,0) node[left] {$B$};
        \draw (nand.out) -- ++(1,0) node[right] {$Y_1=(A\cdot B)'$};
        % Lower circuit: NOT, NOT, OR
        \draw (0,-0.5) node[not port, scale=0.7] (n1) {};
        \draw (0,-2) node[not port, scale=0.7] (n2) {};
        \draw (3,-1.25) node[or port] (or) {};
        \draw (n1.in) -- ++(-1,0) node[left] {$A$};
        \draw (n2.in) -- ++(-1,0) node[left] {$B$};
        \draw (n1.out) -- ++(0.3,0) |- (or.in 1);
        \draw (n2.out) -- ++(0.3,0) |- (or.in 2);
        \draw (or.out) -- ++(1,0) node[right] {$Y_2=A'+B'$};
    \end{circuitikz}
    \caption{The two sides of De Morgan's theorem.}
    \label{fig:demorgan}
\end{figure}

\begin{enumerate}
    \item Build two switch input circuits for $A$ and $B$.
    \item Build both circuits of Figure \ref{fig:demorgan} on the breadboard using the same switches, each output with its own LED indicator stage: $Y_1$ with one NAND gate of the 74LS00, and $Y_2$ with two inverters of the 74LS04 and one OR gate of the 74LS32. Connect Pin 14 and Pin 7 of every IC first.
    \item Complete Table \ref{tab:demorgan}. Is $Y_1 = Y_2$ for all four input combinations? Is $Y_1$ the complement of $A\cdot B$ (compare with the AND truth table)?
    \item Use the multimeter to measure the voltage at the output $Y_1$ for logic 0 and logic 1: \underline{\hspace{1.5cm}} V and \underline{\hspace{1.5cm}} V.
    \item An input of a NAND gate is left unconnected (``floating''). What logic level does the gate read? Why should unused inputs always be tied to a defined level?
\end{enumerate}

\begin{table}[H]
    \centering
    \begin{tabular}{|c|c|c|c|}
        \hline
        $A$ & $B$ & $Y_1=(A\cdot B)'$ & $Y_2=A'+B'$ \\
        \hline
        0 & 0 & & \\
        0 & 1 & & \\
        1 & 0 & & \\
        1 & 1 & & \\
        \hline
    \end{tabular}
    \caption{Verification of De Morgan's theorem.}
    \label{tab:demorgan}
\end{table}

\section*{Part II: Three-Way Light Control in Quartus (about 40 min)}

\begin{enumerate}
    \item Create a new Quartus project \texttt{lab1} for the DE0-CV board and draw the schematic (Block Editor) of the SOP expression of your pre-lab using NOT, AND3 and OR gates. How many gates of each type did you use? \underline{\hspace{4cm}} (You may first draw and test the same circuit in Logisim; the board test still needs Quartus.)
    \item Compile the schematic and perform a functional simulation using a vector waveform file (*.vwf) as in Lab 5.5. Apply all 8 input combinations in order $000, 001, \ldots, 111$ and capture the waveform (in Logisim: toggle the inputs with the Poke tool, or use the chronogram of Logisim-evolution, and capture a screenshot). Compare it with your truth table.
    \item Assign $x_1, x_2, x_3$ to switches SW2, SW1, SW0 and $f$ to LEDR0. Download the circuit to the board and test it: starting from all switches OFF, flip \emph{one} switch at a time and check that the light toggles every time.
\end{enumerate}

\section*{Part III: Car Safety Alarm in Quartus (about 25 min)}

\begin{enumerate}
    \item In the same project, draw \emph{both} circuits $A_1$ and $A_2$ of your pre-lab with the same four inputs $K, D, S, B$.
    \item Simulate (Quartus waveform or Logisim) with all 16 combinations and verify that the two outputs are always identical. Gates in $A_1$: \underline{\hspace{1.5cm}} \; Gates in $A_2$: \underline{\hspace{1.5cm}}
    \item Implement $A_2$ on the board with $K$ = SW0, $D$ = SW1, $S$ = SW2, $B$ = SW3 and $A$ = LEDR0, and test the four scenarios of Table \ref{tab:alarm}.
\end{enumerate}

\begin{table}[H]
    \centering
    \begin{tabular}{|l|c|c|c|c|c|}
        \hline
        \textbf{Scenario} & $K$ & $D$ & $S$ & $B$ & $A$ (LEDR0) \\
        \hline
        Key out, door open & 0 & 0 & 1 & 0 & \\
        Key in, door open & 1 & 0 & 0 & 0 & \\
        Key in, door closed, seated, belt off & 1 & 1 & 1 & 0 & \\
        Key in, door closed, seated, belt on & 1 & 1 & 1 & 1 & \\
        \hline
    \end{tabular}
    \caption{Alarm test scenarios.}
    \label{tab:alarm}
\end{table}

\section*{Extension (take-home, or if time permits)}
\begin{enumerate}
    \item Verify the second De Morgan law, $(A+B)' = A'\cdot B'$, on the breadboard (74LS02 NOR, 74LS04, 74LS08). \textbf{Note that the pin order of the 74LS02 is different:} gate inputs are pins 2,3 $\rightarrow$ 1; 5,6 $\rightarrow$ 4; 8,9 $\rightarrow$ 10; 11,12 $\rightarrow$ 13.
    \item Replace your three-way light circuit with $f = x_1 \oplus x_2 \oplus x_3$ (XOR gates). Show by simulation that it has the same waveform as the SOP circuit. Which one has fewer gates?
    \item The multiplexer is another important example of Lecture 2. You will build and test one on the breadboard in Lab 5.5.
\end{enumerate}

\section*{Check-off}
Show your instructor or teaching assistant:
\begin{enumerate}
    \item Table \ref{tab:demorgan} completed from your breadboard measurements.
    \item The waveform and the working board for the three-way light control.
    \item The equivalence waveform of the two alarm circuits and the working board.
\end{enumerate}

\vspace{1cm}
\noindent \textbf{Updated By:} R. Sutthaweekul \\
\textbf{Version:} 1.0 \\
\textbf{Release Date:} 2026-11-11

\end{document}
